\documentclass[10pt,a4paper]{article}
\usepackage[hmargin=1.8cm,vmargin=1.5cm,includefoot]{geometry}
\usepackage[T1]{fontenc}
\usepackage[utf8]{inputenc}
\usepackage{lmodern}
\usepackage{microtype}
\usepackage{amsmath}
\usepackage{booktabs}
\usepackage{graphicx}
\usepackage{caption}
\usepackage{xcolor}
\usepackage{tcolorbox}
\usepackage{listings}
\usepackage{enumitem}
\usepackage{placeins}
\usepackage{float}
\usepackage[hidelinks]{hyperref}

\captionsetup{font=small,labelfont=bf,skip=4pt}
\setlist{nosep,leftmargin=*}
\setlength{\parskip}{3pt}
\setlength{\parindent}{0pt}
\makeatletter
\setlength{\@fptop}{0pt}
\setlength{\@fpsep}{12pt}
\setlength{\@fpbot}{0pt plus 1fil}
\makeatother
\newcommand{\figdir}{../output/figures}

\lstset{basicstyle=\ttfamily\scriptsize,breaklines=true,columns=fullflexible,frame=single,
  rulecolor=\color{gray!40},keywordstyle=\color{blue!60!black},commentstyle=\color{gray},
  language=Python,showstringspaces=false}

\title{\vspace{-1.2cm}\textbf{How predictable is tomorrow's Bitcoin volatility?}\\[2pt]
\large GARCH, realized variance and an honest 2025-2026 forecast test}
\author{Ignacio Queipo de Llano P\'erez-Gasc\'on\\\small Predictive Statistics in Finance, IESE Master in Finance}
\date{October 2026}

\begin{document}
\maketitle
\vspace{-0.6cm}

\begin{tcolorbox}[colback=white,colframe=black,boxrule=0.4pt,sharp corners,left=6pt,right=6pt,top=4pt,bottom=4pt]
\textbf{Summary.} Bitcoin's direction is not predictable, but its risk is, partly. Volatility clusters strongly, and a
Student-$t$ GARCH describes it far better than a Normal one. Out of sample, over the 608 days from January 2025 to
August 2026, the most important choice is not GARCH versus a simple regression on realized variance, but how much history
the model uses. Bitcoin's volatility today is a fraction of what it was in 2017-2018, so models fitted on all
history forecast risk 30-80\% too high. Re-estimated on the last two years, the simple realized-variance regression is the
best-calibrated and most accurate rule. Even so, it explains under a fifth of the day-to-day variation in risk.
\end{tcolorbox}

\section{Data, timing and how forecasts are judged}

Daily returns are $r_t = 100\,[\log P_t - \log P_{t-1}]$ from 18 August 2017 to 31 August 2026: 3,301 calendar days,
since Bitcoin trades every day. Models are estimated on data up to 31 December 2024 (2,693 days), and every forecast
for day $t+1$ uses only information up to day $t$. The evaluation period is 1 January 2025 to 31 August 2026: 608
one-day-ahead forecasts.

What GARCH forecasts, $\sigma^2_{t+1|t}=\mathrm{Var}(r_{t+1}\mid\mathcal F_t)$, is never observed. I judge it against
realized variance, $RV_{t+1}=\sum_{j=1}^{24} r_{t+1,j}^2$, the sum of the 24 squared hourly returns of that day. RV is
far more informative than one squared daily return, but it is still a noisy proxy. All 608 evaluation days and their
origins have a valid 24-hour RV. Two losses are used. $\text{MSE}=\frac1T\sum(RV_t-\hat h_t)^2$ punishes large
errors in absolute terms, so a few spike days dominate it. $\text{QLIKE}=\frac1T\sum(\log\hat h_t+RV_t/\hat h_t)$ judges
errors relative to the forecast and punishes under-prediction more than over-prediction. Both rank forecasts correctly
even when RV is a noisy proxy (Patton, 2011), which is why the case uses them. Unless stated otherwise, standard errors
are robust: QMLE for GARCH, Newey-West with 7 lags (one week) for regressions.

\section{Part I: is there volatility to forecast?}

\begin{table}[h]
\centering\small
\begin{tabular}{lr@{\hspace{1.4cm}}lr}
\toprule
Mean daily return & 0.11\% & Skewness & $-0.75$\\
Standard deviation & 3.71\% (71\% a year) & Excess kurtosis & 12.5\\
Worst / best day & $-46.5\%$ / $+22.5\%$ & Days beyond 4 s.d. (Normal: 0.17) & 12\\
\midrule
Ljung-Box, $r_t$, 10 / 20 lags & $p=0.031$ / $0.076$ & Ljung-Box, $r_t^2$, 10 / 20 lags & $p<10^{-15}$\\
Largest $|\text{ACF}(r_t)|$, lags 1-20 & 0.045 & ACF$(r_t^2)$ at lag 1 & 0.099\\
ARCH-LM, 1 / 5 / 10 lags & \multicolumn{3}{l}{$p=4\times10^{-7}$ / $5\times10^{-9}$ / $7\times10^{-12}$}\\
\bottomrule
\end{tabular}
\caption{Training sample, 18 Aug 2017 to 31 Dec 2024 (2,693 days). The 95\% band for one autocorrelation is $\pm0.038$.}
\label{tab:t1}
\end{table}

\begin{figure}[h]
\centering
\includegraphics[width=0.7\linewidth]{\figdir/f2_acf.png}
\caption{Autocorrelations in the training sample. Returns stay inside the noise band; squared returns stay above it.}
\label{fig:acf}
\end{figure}

\textbf{(a) Returns are not meaningfully predictable from their own past.} No return autocorrelation exceeds 0.045 in
absolute value, which means yesterday's return explains at most 0.2\% of today's. The Ljung-Box test rejects at 10 lags at
the 5\% level but not at 20. With 2,693 observations that is weak evidence of a tiny effect, not a usable
signal.

\textbf{(b) Volatility clusters.} Squared returns are positively autocorrelated, and the Ljung-Box statistic for
$r_t^2$ is five times that for $r_t$. Large moves, of either sign, tend to follow large moves. The distribution
is also far from Normal: twelve days beyond four standard deviations where a Normal would give 0.17, led by
$-46.5\%$ on 12 March 2020.

\textbf{(c) What the ARCH-LM test does and does not establish.} Rejecting the null hypothesis that squared residuals
are uncorrelated establishes that the conditional variance is not constant. It does not establish that a GARCH model
is the right description, how large or persistent the effect is, or that forecasts will be accurate. A change in the
level of volatility between periods also produces autocorrelated squared returns without day-to-day clustering, and
Bitcoin had such changes (Part VI). The test also says nothing about the mean: predictable risk does not mean
predictable returns.

\section{Part II and III: what GARCH learns, and why the distribution matters}

\begin{table}[h]
\centering\small
\begin{tabular}{lrrr}
\toprule
& ARCH(1) Normal & GARCH(1,1) Normal & GARCH(1,1) Student-$t$\\
\midrule
$\omega$ & 11.90 (9.3) & 0.66 (2.4) & 0.13 (1.5)\\
$\alpha$: reaction to yesterday's shock & 0.136 (3.0) & 0.100 (2.4) & 0.073 (6.5)\\
$\beta$: weight on yesterday's variance & -- & 0.859 (19.2) & 0.927 (61.2)\\
$\alpha+\beta$ [95\% CI] & 0.136 & 0.958 [0.925, 0.992] & 1.000 [0.984, 1.016]\\
Half-life of a volatility shock & 0.3 days & 16 days & no decay\\
$\nu$ (degrees of freedom) & -- & -- & 3.18 (20.1)\\
Log-likelihood & $-7{,}307$ & $-7{,}165$ & $-6{,}831$\\
AIC / BIC & 14,620 / 14,638 & 14,337 / 14,361 & 13,673 / 13,702\\
\bottomrule
\end{tabular}
\caption{Estimates on the training sample, constant mean, robust $t$-statistics in brackets. The interval for
$\alpha+\beta$ uses the joint covariance: $\hat\alpha$ and $\hat\beta$ are strongly negatively correlated ($-0.92$ and $-0.85$).}
\label{tab:garch}
\end{table}

\textbf{Why GARCH helps.} ARCH(1) lets only yesterday's shock move today's variance, so a shock is forgotten within a
day (half-life 0.3 days). This contradicts Figure~\ref{fig:acf}, where squared returns stay correlated for weeks. Adding
yesterday's variance ($\beta$) gives the model memory: GARCH(1,1) improves the log-likelihood by 142 points for one
extra parameter, and both AIC and BIC prefer it clearly.

\textbf{Reaction and persistence.} In the Normal GARCH, a 10\% daily move adds about $0.10\times10^2\approx10$
to tomorrow's variance (in \%$^2$). Starting from a calm 2.5\% daily volatility, tomorrow's volatility jumps to about
4\%. The excess then decays at rate $\alpha+\beta=0.958$ a day: half of it is still there after 16 days. In the
Student-$t$ GARCH the immediate reaction is smaller ($\alpha=0.073$), but the estimated persistence reaches the
boundary of 1: shocks never fade and there is no long-run variance to return to. This is an integrated GARCH,
equivalent to an exponentially weighted moving average of squared returns. Part VI suggests why it appears.

\textbf{Significance is not forecast accuracy.} $\alpha$ and $\beta$ are clearly significant, but a significant
$t$-statistic only says the parameter is not zero \emph{inside the sample, if the model is right}. Forecast accuracy
depends on other things too. The parameters must stay stable after the sample ends. The model's implied variance level
must match the future level. And RV must not be so noisy that the gains disappear. With 2,693 days even modest effects are
significant. Here $\alpha$ and $\beta$ are individually blurred by their correlation, even though their sum is
estimated precisely.

\textbf{Normal versus Student-$t$.} The $t$ specification improves the log-likelihood by 333 points for one extra parameter
($\nu$). $\hat\nu=3.18$ (s.e.\ 0.16) means very heavy tails: with $\nu\le4$ the fourth moment does not exist, so
kurtosis is infinite. Volatility clustering does not explain the tails away. The standardized residuals of the Normal
GARCH still have excess kurtosis of 13.5. It helps to keep two choices apart. The \emph{variance recursion}
$\sigma^2_t=\omega+\alpha\varepsilon^2_{t-1}+\beta\sigma^2_{t-1}$ is the same in both models; only the
\emph{distribution of $z_t$} changes. But that choice still changes the recursion's estimates. Under a Normal, a 15\%
day is almost impossible, so the likelihood raises $\alpha$ and $\omega$ to make it less surprising. Under a $t$,
the same day is an ordinary tail draw, so variance reacts less to each shock and persists more. For risk measures the
distribution also matters directly: the same $\sigma$ implies a much fatter left tail, which will matter for VaR and
Expected Shortfall.

\section{Part IV and V: forecasting a latent variable, and a simple benchmark}

Each day from 1 January 2025 I re-estimate GARCH-$t$ on all returns up to day $t$ (expanding window) and forecast
$\hat h_{t+1}=\hat\sigma^2_{t+1|t}$. The colleague's benchmark, \emph{RV-OLS}, regresses tomorrow's RV on today's:
$RV_{t+1}=a+b\,RV_t+u_{t+1}$, re-estimated in the same way on the same dates.\footnote{Re-estimating RV-OLS at every
origin makes it comparable to GARCH and to Part VI. Frozen at the end-2024 estimates, it does slightly worse: MSE
62.7, QLIKE 2.782.}

\textbf{The benchmark in sample.} On the training data $\hat b=0.548$ (Newey-West $t=12.5$; plain OLS would claim
33), with 95\% interval [0.46, 0.63] and $R^2=0.29$. Today's RV clearly carries significant information about tomorrow's.
But $b=0.55$ means a jump in RV loses half its effect within about a day: the forecast quickly returns to the
regression's long-run mean, $\hat a/(1-\hat b)=16.2$.

\begin{table}[h]
\centering\small
\begin{tabular}{lrrrrrr}
\toprule
Rule & MSE & QLIKE & Corr$(\hat h,RV)$ & Mean $\hat h$ & $\hat h$ / mean RV & MZ slope\\
\midrule
GARCH-$t$, expanding & 48.7 & 2.612 & 0.33 & 7.12 & 1.32 & 0.53\\
RV-OLS, expanding & 58.1 & 2.753 & \textbf{0.43} & 9.76 & 1.82 & 0.77\\
GARCH-$t$, rolling 2 years & 47.0 & 2.615 & 0.29 & 6.62 & 1.23 & 0.54\\
RV-OLS, rolling 2 years & \textbf{39.0} & \textbf{2.590} & 0.42 & 5.97 & \textbf{1.11} & 1.22\\
\bottomrule
\end{tabular}
\caption{Out-of-sample scores, 608 days, 1 Jan 2025 to 31 Aug 2026. Mean realized variance: 5.38 (2.3\% daily
volatility). MZ slope: $b$ in the Mincer-Zarnowitz regression $RV_t=a+b\,\hat h_t$; a well-scaled forecast has $b=1$.
The rolling rows belong to Part VI and are shown here for comparison.}
\label{tab:oos}
\end{table}

\begin{figure}[t]
\centering
\includegraphics[width=0.76\linewidth]{\figdir/f3_forecasts.png}
\caption{One-day-ahead variance forecasts against realized variance, 2025-2026. Spikes are cut at 60.}
\label{fig:fc}
\end{figure}

\textbf{What each statistic says.} MSE measures accuracy in absolute terms and is driven by a handful of spike days:
half-yearly MSE ranges from 16 to 86. QLIKE measures errors relative to the forecast's own level, so a forecast of 2
when RV turns out to be 10 costs much more than the reverse. Correlation measures only whether the forecast moves
\emph{with} RV. It ignores the level, so a forecast that is consistently 80\% too high keeps its correlation and still
has poor MSE and QLIKE. That is RV-OLS expanding: it has the highest correlation (0.43) and the worst losses, because
it forecasts 9.8 on average when realized variance averages 5.4, and falls short of RV on only 12\% of days. GARCH-$t$
is also miscalibrated, 32\% too high on average. Its MZ slope of 0.53 says it is also too reactive: when it forecasts high,
RV tends to come in lower.

\textbf{A quiet day does not prove a high forecast wrong.} The forecast is the variance of a distribution, not a
prediction of one number. If tomorrow's true variance is high, a calm day is still quite likely, especially with
$\nu\approx3$, where most draws of $z^2$ are small and a few are huge. RV is in turn a noisy reading of the true
variance. One observation cannot separate a bad forecast from an ordinary draw. Calibration can only be judged over
many days, as in Table~\ref{tab:oos}.

\textbf{Significance versus usefulness.} Both questions get a different answer. Does today's RV carry significant
information about tomorrow's? Yes, with $t=12.5$. Does the model forecast new days well enough to be useful? Only
modestly: the best rule's Mincer-Zarnowitz $R^2$ is 0.18, so more than four fifths of the daily variation in RV is
unforecastable noise. A coefficient can be very precisely different from zero and still leave most of the variation
unexplained. Its intercept can also come from a period that no longer applies, as here.

\textbf{GARCH versus RV-OLS dynamics} (Figure~\ref{fig:zoom}). The two rules react to different things. RV-OLS is
almost a rescaled copy of yesterday's RV (correlation 0.997): it jumps after every intraday storm and forgets it within
two days (average daily change 2.3). GARCH only sees the \emph{daily} return. It moves smoothly (average daily change
0.5), reacts to large closing moves and keeps them in memory for weeks. On 20 January 2025, RV hit 60 while the day closed
only +0.9\%: GARCH did not react at all, and RV-OLS jumped to 23 for a day that came in at 12. On 5-6 February
2026, returns of $-15.2\%$ and $+11.8\%$ took GARCH from 9 to 40, and a week later it still forecast 21 while RV was back
at 5-10. RV-OLS was back at 7 within two days.

\begin{figure}[h]
\centering
\includegraphics[width=0.82\linewidth]{\figdir/f4_zoom_spike.png}
\caption{Two episodes, rolling-window forecasts. GARCH misses intraday storms with flat closes and holds onto large
closing moves for weeks; RV-OLS reacts to everything for a day.}
\label{fig:zoom}
\end{figure}

GARCH is more dynamically flexible in one sense: its memory depends on the state, and it builds up and fades
gradually. In 2025-2026, however, volatility spikes were short-lived, so that memory mostly produced over-prediction
after shocks. With expanding windows GARCH-$t$ beats RV-OLS on QLIKE (Diebold-Mariano $t=-5.6$, $p<0.001$). Part VI
shows that this advantage comes from calibration, not from the dynamics.

\section{Part VI: does distant Bitcoin history still help?}

I repeat the exercise with a two-year rolling window: the forecast for 1 January 2025 uses 1 January 2023 to 31
December 2024, and the window moves one day at a time. The evaluation dates do not change.

\textbf{(a) Rolling versus expanding.} For RV-OLS, the two-year window cuts MSE by a third (58.1 to 39.0) and QLIKE from
2.753 to 2.590, both strongly significant (Diebold-Mariano $t\approx-8$). Its average forecast falls from 1.82 to 1.11
times realized variance, while its correlation does not change. The gain is entirely in the level. For GARCH-$t$ the
window barely matters (MSE 47.0 vs 48.7, QLIKE 2.615 vs 2.612; neither difference significant). With the rolling window
the ranking flips: RV-OLS beats GARCH-$t$ on both losses, although only at the 10\% level (MSE $p=0.07$, QLIKE $p=0.10$).

\begin{figure}[h]
\centering
\includegraphics[width=0.85\linewidth]{\figdir/f5_rolling_params.png}
\caption{Rolling two-year estimates at each forecast origin. Left: RV-OLS slope with Newey-West 95\% band. Right:
GARCH-$t$ parameters.}
\label{fig:roll}
\end{figure}

\textbf{(b) The RV persistence coefficient} is positive and significant at every one of the 608 origins, but lower
than in the full sample. It rises slowly from 0.30 to 0.42, against 0.55 with all history, and the rolling $R^2$ is
0.09-0.17 rather than 0.29. Today's RV is reliably informative, but less so in the recent, calmer period.

\textbf{(c) GARCH persistence.} With all history, $\alpha+\beta$ sits at 1.000 at every origin. With two years, it is
1.000 until September 2025 and then drifts between 0.95 and 0.98. The 95\% interval contains 1 throughout, while
$\alpha$ ranges from 0.02 to 0.09 and $\beta$ moves to compensate. Persistence close to one means shocks never die out:
the model has no long-run variance to revert to, and its multi-day forecasts stay flat at today's level. A well-known
cause is unmodelled shifts in the level of variance (Lamoureux and Lastrapes, 1990). A GARCH with one long-run level,
fitted to a series whose level drifts down, explains the drift as extremely persistent shocks. The first rolling windows
include the rise from 2023 to 2024 (mean RV 4.9 to 7.6), and persistence leaves the boundary only once that rise drops
out of the window.

\textbf{(d) What the ranking change says about distant data.} Mean RV was 58.6 in 2017, 25.8 in 2018 and has been
about 5 since 2023 (Figure~\ref{fig:regimes}, appendix). Old observations describe a different level of risk. The rule
hurt most is the one that leans hardest on an estimated long-run mean. RV-OLS forgets within a day, so on most days its
forecast \emph{is} its long-run mean, and with all history that mean is 16 against a realized 5. GARCH-$t$ with
$\alpha+\beta\approx1$ barely uses a long-run mean, so it suffers less from old data. That gave it its expanding-window
advantage. Distant history therefore does not help: it mostly contaminates the level.

\textbf{(e) This does not prove a structural break.} The same pattern would come from a gradual decline in volatility,
from switching between regimes, or from long memory in volatility. Realized variance is known to have long memory,
which is why models that combine daily, weekly and monthly averages of RV work well. I did not test for a break, and the
evidence comes from a single 20-month path in which a few spikes decide the MSE. Part of the gap is also not economic.
In 2017-2018 the sum of hourly squared returns was 1.4-1.7 times the squared daily returns. Since 2019 the ratio has
been about 1. The early RV proxy looks inflated, probably by noisier intraday prices in thinner markets, so old data are
partly worse data, not just data from another era.

\section{Part VII: recommendation to the risk manager}

\begin{itemize}
\item \textbf{Is there a predictable component?} In risk, yes; in direction, no. Today's volatility reliably signals
tomorrow's, but at best about 18\% of daily variation in realized variance is forecastable.
\item \textbf{Is GARCH economically sensible?} Yes: it captures clustering, and the Student-$t$ version
($\nu\approx3$) is clearly better for Bitcoin. Its persistence estimate at one is a warning, not a feature.
\item \textbf{Forecast quality against RV.} All rules over-forecast risk. The model ranking depends on the window, but
the window effect itself is large, significant and consistent.
\item \textbf{Significance versus out-of-sample.} Every key parameter is significant, and that told us nothing about
which rule would forecast best.
\item \textbf{Information set.} Rolling two-year windows. The cost is noisier parameters (Figure~\ref{fig:roll}); the
benefit is a forecast level that matches the current regime.
\item \textbf{Uses.} Volatility targeting and position sizing, risk limits, margin buffers, and as the $\sigma$ input
to VaR and Expected Shortfall (with Student-$t$ tails). Not as a precise forecast of any single day.
\item \textbf{Main uncertainty.} One 20-month test period, dominated by a few spikes. The GARCH-versus-RV ranking is
only marginally significant and could reverse in a more persistent volatile period, like 2022, where GARCH's
memory would pay off.
\end{itemize}

\begin{tcolorbox}[colback=white,colframe=black,boxrule=0.4pt,sharp corners,left=6pt,right=6pt,top=4pt,bottom=4pt]
\textbf{Executive recommendation.} Tomorrow I would implement the realized-variance regression, re-estimated daily on
the last two years. It was the most accurate and best-calibrated rule over 2025-2026, and it is simple enough to audit.
GARCH-$t$ on the same window would run alongside as a second opinion after large closing moves, where it reacts more
lastingly. I would use the forecast to scale Bitcoin positions to a volatility target, to set daily risk limits and margin
buffers, and as the volatility input to VaR and Expected Shortfall with fat tails. Confidence should be moderate on the
\emph{level} of risk and low on any single day: four fifths of daily variation is unpredictable, and both models still
over-forecast risk on average. I would review the window and the calibration every quarter, because the most important lesson
from this exercise is that Bitcoin's risk level changes and old data mislead.
\end{tcolorbox}

\small
\textbf{References.} Lamoureux, C. and Lastrapes, W. (1990), Persistence in variance, structural change and the GARCH
model, \emph{JBES} 8(2). Patton, A. (2011), Volatility forecast comparison using imperfect volatility proxies,
\emph{Journal of Econometrics} 160(1).
\normalsize

\clearpage
\appendix
\section{Additional figures}

\begin{figure}[H]
\centering
\includegraphics[width=0.9\linewidth]{\figdir/f6_regimes.png}
\caption{Average daily realized variance and squared daily return by year. The shaded years are the evaluation
period. Until 2018, RV is well above the squared daily return; since 2019 the two agree.}
\label{fig:regimes}
\end{figure}

\begin{figure}[H]
\centering
\includegraphics[width=0.9\linewidth]{\figdir/f7_cum_qlike.png}
\caption{Cumulative QLIKE of RV-OLS minus GARCH-$t$; rising means GARCH-$t$ is ahead. With expanding windows GARCH
pulls ahead steadily from mid-2025. With rolling windows RV-OLS stays slightly ahead.}
\label{fig:cum}
\end{figure}

\begin{figure}[H]
\centering
\includegraphics[width=0.9\linewidth]{\figdir/f1_price_returns.png}
\caption{Bitcoin price and daily returns, 2017-2026. The dashed line marks the start of the evaluation period.}
\label{fig:price}
\end{figure}

\section{Python code}

The full code is in the repository (\texttt{src/}): \texttt{task1.py} (Part I), \texttt{task2\_3.py} (Parts II-III),
\texttt{forecast.py} (all out-of-sample forecasts), \texttt{evaluate.py} (Parts IV-VI) and \texttt{figures.py}.
Python 3.13 with pandas, statsmodels and \texttt{arch}. The core of the exercise, the forecast loop, is reproduced
below.

\lstinputlisting{../src/forecast.py}

\end{document}
